% =====================================================================
%  1bat-ud5-14.tex  —  SESSIÓ 14: Repàs global i síntesi (estacions)
%  (sense preàmbul: s'inclou des de main.tex)
% =====================================================================
\horatitol{Sessió 14 — Repàs global i síntesi (estacions)}%
{Repàs actiu de tota la unitat per quatre estacions, una per bloc, i un mapa final que connecta cada eina amb la realitat econòmica i social que descriu.}

\subsection{Idea clau}

Avui repassem \textbf{tota la unitat} de manera activa. Hi ha \textbf{quatre
estacions}, una per cada bloc treballat. Rotareu per totes amb temps controlat (uns
$10$ minuts cada una), de manera que tornareu a tocar tots els continguts en contextos
diferents. Al final, construirem entre tots un \textbf{mapa de la unitat}.

\begin{idea}
\textbf{Les quatre estacions}
\begin{enumerate}[leftmargin=1.6em, itemsep=2pt]
  \item \textbf{Percentatges:} factor multiplicador, augments/disminucions i
        encadenats (multiplicar índexs).
  \item \textbf{Interès i crèdits:} interès compost, TIN i TAE, comparar ofertes.
  \item \textbf{IPC i poder adquisitiu:} actualitzar una quota i calcular la pèrdua de
        poder de compra.
  \item \textbf{Pressupost:} estructurar despeses i ingressos i calcular el balanç.
\end{enumerate}
\textbf{Consell:} a cada estació, abans de calcular, pregunta't \emph{quina} realitat
econòmica o social hi ha al darrere.
\end{idea}

\subsection{Les estacions}

\begin{exemple}[Estació 1 — percentatges encadenats]
Un material de $\num{150}\eur$ puja un $20\%$ i, sobre el preu resultant, té un
descompte del $15\%$. Quin és el preu final? A quina variació global equival?
\end{exemple}

\begin{exemple}[Estació 2 — interès i TAE]
Un microcrèdit ofereix un TIN del $4\%$ mensual. Calcula'n la TAE. Et sembla un crèdit
barat o car? Què aconsellaries a la família?
\end{exemple}

\subsection{Exercicis resolts}

\begin{exresolt}[model de l'Estació 1: encadenats]
Per al material de $\num{150}\eur$ (puja $20\%$, després baixa $15\%$), calcula el preu
final i la variació global.
\solucio
Multipliquem els índexs:
\[
\num{150}\per 1,20\per 0,85=\num{150}\per 1,02=\num{153}\eur.
\]
L'índex global és $1,02$, és a dir una \textbf{pujada global del $2\%$} (no del $5\%$
de restar $20-15$, ni cap compensació exacta).
\resposta{$\num{153}\eur$; pujada global del $2\%$ (índex $1,02$).}
\end{exresolt}

\begin{exresolt}[model de l'Estació 2: TAE]
Per a un TIN del $4\%$ mensual, calcula la TAE i comenta-la.
\solucio
Si cada mes es multiplica per $1,04$, en un any es multiplica per $1,04^{12}$:
\[
\text{TAE}=(1,04)^{12}-1\approx 1,601-1=0,601=60,1\%.
\]
Un «$4\%$» mensual és una TAE del $60,1\%$: \textbf{caríssim}. Aconsellaria a la
família \textbf{no} acceptar-lo i buscar alternatives (préstec personal de banc, ajuts,
serveis socials) amb una TAE molt més baixa.
\resposta{TAE $\approx 60,1\%$; és un crèdit molt car, desaconsellable.}
\end{exresolt}

\subsection{Exercicis per fer}

\noindent\textit{Una tasca per estació. Rota i resol-les totes.}

\begin{exercicis}
\begin{exlist}
  \item \textbf{(Estació 1 — percentatges)} Una factura de $\num{200}\eur$ puja un
        $10\%$ i després baixa un $5\%$. Calcula el preu final i la variació global.
  \item \textbf{(Estació 2 — interès i TAE)} Calcula el capital final de $\num{1000}\eur$
        al $3\%$ compost anual durant $3$ anys. Després, calcula la TAE d'un TIN del
        $5\%$ mensual.
  \item \textbf{(Estació 3 — IPC)} Un lloguer de $\num{700}\eur$ s'actualitza amb un IPC
        del $3,5\%$. Calcula la nova quota. Si l'ajut de la família ($\num{950}\eur$) es
        congela, quant li queda ara per a la resta de despeses?
  \item \textbf{(Estació 4 — pressupost)} Un projecte té una subvenció de $\num{1200}\eur$
        i $20$ quotes de $\num{35}\eur$; les despeses sumen $\num{1750}\eur$. Calcula
        els ingressos i el balanç. Quadra?
  \item \textbf{(Mapa de la unitat)} Completa, amb les teves paraules, la connexió de
        cada eina amb la realitat que descriu:
        \begin{center}
        \renewcommand{\arraystretch}{1.4}
        \begin{tabular}{p{4.8cm} p{7.8cm}}
        \toprule
        \textbf{Eina} & \textbf{Quina realitat descriu / per a què serveix} \\
        \midrule
        Factor i índex de variació & \rule{7.2cm}{0.4pt} \\
        Percentatges encadenats & \rule{7.2cm}{0.4pt} \\
        Interès compost / TAE & \rule{7.2cm}{0.4pt} \\
        IPC i poder adquisitiu & \rule{7.2cm}{0.4pt} \\
        Pressupost i balanç & \rule{7.2cm}{0.4pt} \\
        \bottomrule
        \end{tabular}
        \end{center}
  \item \textbf{(Síntesi)} En una frase per bloc, resumeix què has après a la unitat:
        calcular variacions percentuals, entendre el cost d'un crèdit, interpretar la
        inflació i quadrar un pressupost.
\end{exlist}
\end{exercicis}

% ---------------------------------------------------------------------
\fullsolucions{Sessió 14}
\begin{solucions}
\begin{sollist}
  \item $\num{200}\per 1,10\per 0,95=\num{200}\per 1,045=\num{209}\eur$. Índex global
        $1,045$: pujada global del $4,5\%$.
  \item Capital final: $\num{1000}\per (1,03)^3\approx\num{1092,73}\eur$. TAE d'un $5\%$
        mensual: $(1,05)^{12}-1\approx 0,7959=79,6\%$.
  \item Nova quota: $\num{700}\per 1,035=\num{724,5}\eur$. Amb l'ajut congelat de
        $\num{950}\eur$, ara li queden $\num{950}-\num{724,5}=\num{225,5}\eur$ (abans li
        quedaven $\num{950}-\num{700}=\num{250}\eur$: ha perdut $\num{24,5}\eur$ al mes).
  \item Ingressos: $\num{1200}+20\per\num{35}=\num{1200}+\num{700}=\num{1900}\eur$.
        Balanç: $\num{1900}-\num{1750}=\num{150}\eur$. És positiu: el projecte
        \textbf{quadra}.
  \item Resposta oberta. Idees: \emph{factor/índex} $\rightarrow$ aplicar pujades i
        baixades d'un sol cop; \emph{encadenats} $\rightarrow$ combinar diversos canvis
        (multiplicant); \emph{interès compost/TAE} $\rightarrow$ veure el cost real d'un
        crèdit i detectar abusos; \emph{IPC i poder adquisitiu} $\rightarrow$ entendre
        com la inflació afecta famílies i ajuts; \emph{pressupost i balanç}
        $\rightarrow$ planificar i justificar econòmicament un projecte.
  \item Resposta oberta. Exemple: «He après a \emph{calcular} variacions percentuals amb
        índexs, a \emph{entendre} per què un crèdit ràpid és abusiu (TAE), a
        \emph{interpretar} la pèrdua de poder de compra amb la inflació i a
        \emph{quadrar} el pressupost d'un projecte amb el full de càlcul.»
\end{sollist}
\end{solucions}
